\section{Flexible Bayesian Prior}
We want to assess if the additional flexibility of its skewness and kurtosis figures improves the accuracy of Bayesian estimation, by employing it as a highly flexible Bayes prior. We will also assess the computational costs that come with the estimation of 2 additional parameters, especially when they rely on numerical computation, and whether this is an acceptable compromise.

We create a simple Bayesian linear regression setup based solely on the \sindist{} given its great approximation capabilities, see Mathematical properties as well as empirical results.
\begin{align*}
	y_i &= \beta_0 + \vec x_i \beta_1 \epsilon_i \\
	\beta_0 &\sim \Sinu(a_0, d_0, s_0, k_0)\\
	\beta_1 &\sim \Sinu(a_1, d_1, s_1, k_1)\\
	\epsilon_i &\sim \Sinu(-d_Y, d_Y, s_Y, k_Y) \text{ centered at 0}\\
	\vector{\hat \beta} &= (\mat X^\intercal \mat X)^{-1} \mat X^\intercal \vector Y\\
	d_Y &\sim \Sinu(a_d, d_d, s_d, k_d)\\
	s_Y &\sim \Sinu(a_s, d_s, s_s, k_s)\\
	k_Y &\sim \Sinu(a_k, d_k, s_k, k_k)
\end{align*}
There is no hope for conjugacy here. We attempt MCMC via Metropolis-Hastings thereby.

We create an entire linear regression model with every distribution some \sindist{}.

\begin{align*}
	\epsilon_i &\sim \Sinu(a_y,d_y,s_y,k_y)\\
	a &\sim \Sinu(a_a, d_a, s_a, k_a)\\
	d
\end{align*}

